\section{Disjoint-label products and first occurrences}
\label{sec:first}

The first occurrence of a label has a direct monomial description:
fill each row of its Young diagram with the least-weight one-label
monomials.  Each column then forms a lower monomial diagram.
We generate the alternants of these diagrams on disjoint label
sets and apply the Young symmetrizer.  This generates every
minimum-degree filling.

For $V=K^k$, define
\[
 b_j=\dim\Sym^jV=\binom{k+j-1}{k-1},\qquad
 N_j=\binom{k+j}{k},\qquad N_{-1}=0.
\]
Let $d_i=j$ for $N_{j-1}<i\leq N_j$, and pad all partitions with
zeros.  Thus $d_1,d_2,\ldots$ are the degrees of all monomials in $k$
variables, in nondecreasing order.

\begin{theorem}[Complete first occurrences]
\label{thm:r5-complete-first}
For every $n,k\geq1$ and every $\mu\vdash n$, put
\[
 \delta_\mu(k)=\sum_{i\geq1}\mu_i d_i.
\]
There is no $S^\mu$ in $P_{n,k}$ below this degree, and
\begin{equation}\label{eq:r5-first-equality}
 \Pol_{n,k}[\delta_\mu(k)]_\mu
 =\Harm_{n,k}[\delta_\mu(k)]_\mu
 =P_{n,k}[\delta_\mu(k)]_\mu.
\end{equation}
Its row multiplicity space is
\begin{equation}\label{eq:r5-first-character}
 \Hom_{S_n}\bigl(S^\mu,\Pol_{n,k}[\delta_\mu(k)]\bigr)
 \cong
 \bigotimes_{j\geq0}
 \Schur_{(\mu_{N_{j-1}+1},\ldots,\mu_{N_j})}(\Sym^j V).
\end{equation}
\end{theorem}

The degree and tensor-product formula come from the semistandard
construction of Schur functors; Romero records the ambient
plethystic character in \cite[Theorem~1.2 and \S2]{r5-Romero}.
For the alternating label, Wallach gives the determinant-twisted
exterior-power description of the first layer
\cite[Theorem~1]{r5-Wallach}.  The assertion here generates that
entire space.  Its proof rests on a product theorem for arbitrary
lower monomial diagrams.

\subsection{Curvilinear quotients and colored determinants}
\label{sec:r5-quot}

A finite nonempty set $\Gamma$ of monomials in $u_1,\ldots,u_k$ is a
\emph{lower diagram} if it is closed under divisors.
For labels $C$ with $|C|=|\Gamma|$, put
\[
 \Delta_\Gamma(X_C)=
       \det(X_i^\alpha)_{i\in C,\,u^\alpha\in\Gamma},
 \qquad
 X_i^\alpha=\prod_{a=1}^k(x_i^{(a)})^{\alpha_a}.
\]
An ordering of $C$ and $\Gamma$ fixes the irrelevant overall sign.

\begin{theorem}[Products of lower-diagram alternants]
\label{thm:r5-products}
Partition $\{1,\ldots,n\}$ into nonempty sets
$C_1,\ldots,C_r$.  For each $c$, choose any lower diagram
$\Gamma_c$ of $|C_c|$ monomials in $k$ variables.  Then
\begin{equation}\label{eq:r5-product}
       \prod_{c=1}^r\Delta_{\Gamma_c}(X_{C_c})
                         \ \in\ \Pol_{n,k}.
\end{equation}
The lower diagrams may be chosen independently.
The derivative and polarization module generated by this product is
also contained in $\Pol_{n,k}$.
\end{theorem}

The product theorem follows by placing all label blocks in one
quotient of a free module.  B\'erczi--Svendsen's theorem moves that
monomial quotient into the closure of the curvilinear locus, where
a jet determinant gives its operator realization.  We work over
$\C$ with the underlying reduced loci and then descend the resulting
linear-span statement to $K$.  The case $n=1$ is immediate.

For $S=\C[u_1,\ldots,u_k]$, the
punctual Quot scheme parametrizes length-$n$ quotients of $S^{\oplus r}$
supported at the origin.  Its curvilinear locus consists of quotients
whose underlying $S$-module is the structure module of an embedded
curvilinear length-$n$ subscheme.  In coordinates, such a quotient has
the form
\begin{equation}\label{eq:r5-curvilinear-quotient}
 S^{\oplus r}\longrightarrow A=\C[t]/(t^n),\qquad
 u_a\longmapsto f_a(t),\qquad e_c\longmapsto g_c(t),
\end{equation}
where the $f_a$ are pointed jets, at least one has nonzero linear
coefficient, and the $g_c$ generate $A$ as an $S$-module.  Equivalently,
at least one $g_c$ has nonzero constant term.

\begin{theorem}[B\'erczi--Svendsen]
\label{thm:r5-bs-input}
Every quotient
\[
               S/I_1\oplus\cdots\oplus S/I_r
\]
with each $I_c$ a monomial ideal of finite colength is contained in
the closure of the curvilinear locus in the length-$n$ punctual
Quot scheme, where $n=\sum_c\dim_\C S/I_c$.
\end{theorem}

This is \cite[Theorem~1.3, equivalently Theorem~5.2]{r5-BS}.
The same paper's Theorem~1.1 gives the $r=1$ Hilbert-scheme case.
The rank-$r$ Quot statement treats the direct sum of the independent
monomial quotients and is the input for \eqref{eq:r5-product}.

\subsubsection{The graph determinant}

For pointed jets $f=(f_1,\ldots,f_k)$ modulo $t^n$, define
\begin{equation}\label{eq:r5-jet-determinant}
 G_f(X)=\det\left([t^j]\exp\left(\sum_{a=1}^k
                                 x_i^{(a)}f_a(t)\right)\right)
                _{\substack{1\leq i\leq n\\0\leq j<n}}.
\end{equation}
This is the jet determinant of
Theorem~\ref{thm:r3-normal-form}, with the constant
$\prod_{j=0}^{n-1}j!$ omitted. On the graph chart
$f_1(t)=t$, $f_a(t)=\sum_s a_{a,s}t^s$, that theorem gives
\begin{equation}\label{eq:r5-graph-operator}
 \left(\prod_{j=0}^{n-1}j!\right)G_f
 =\exp\left(\sum_{a=2}^k\sum_{s=1}^{n-1}a_{a,s}Y_{a,s}\right)
       \Delta_n\in\Pol_{n,k},
 \qquad Y_{a,s}=\sum_i x_i^{(a)}\partial_{x_i}^{s}.
\end{equation}
Its inverse source reparametrization gives the same membership
throughout the dense regular-jet locus. We now retain this
membership under the colored specialization.

\subsubsection{The colored Pl\"ucker map}

To carry curvilinear closure into a linear span of polynomials,
we use one linear map on a fixed exterior power.  Its label colors
remember which lower diagram belongs to each block.  At a graph
quotient it becomes a differential operator on the jet determinant;
at a monomial quotient it becomes the desired product of alternants.

Every length-$n$ module supported at the origin is annihilated by
$\m^n$, where $\m=(u_1,\ldots,u_k)$.  To see this, the strict
nonzero steps in its $\m$-adic filtration each decrease dimension;
a repeated step would be zero by Nakayama's lemma.  All quotient
points therefore belong to the fixed Grassmannian of $n$-dimensional
quotients of
\[
                     W=(S/\m^n)^{\oplus r}.
\]
Let $c(i)=c$ for $i\in C_c$.  Define a linear map
\begin{equation}\label{eq:r5-colored-map}
 \begin{split}
 L_C:\bigwedge^nW^*&\longrightarrow P_{n,k},\\
 q_0\wedge\cdots\wedge q_{n-1}
   &\longmapsto
 \det\left(\sum_{|\alpha|<n}
        q_j(u^\alpha e_{c(i)})\frac{X_i^\alpha}{\alpha!}\right)_{i,j},
 \qquad \alpha!=\prod_a\alpha_a!.
 \end{split}
\end{equation}
Multilinearity and alternation in the $q_j$ make this a well-defined
linear map on the exterior power.  Its normalization does not depend
on the quotient point.  Its values have total degree at most
$n(n-1)$, so they lie in a fixed finite-dimensional polynomial
space $E=P_{n,k,\le n(n-1)}$.

For the graph quotient \eqref{eq:r5-curvilinear-quotient}, take
$q_j$ to be its coefficient-of-$t^j$ functional.  Then
\begin{align}
 L_C(q_0\wedge\cdots\wedge q_{n-1})
  &=\det\left([t^j]g_{c(i)}(t)
                \exp\left(\sum_a x_i^{(a)}f_a(t)\right)\right)_{i,j}
       \label{eq:r5-curv-image}\\
  &=\left(\prod_{i=1}^n g_{c(i)}(\partial_{x_i})\right)G_f
       \ \in\ \Pol_{n,k}.\nonumber
\end{align}
The second equality uses $f_1(t)=t$.  Each derivative affects just
one row of the determinant.  Constant terms of the $g_c$ are allowed
and act as scalar summands of these differential operators.

The graph chart is dense in the curvilinear Quot locus: the regular
jets form an irreducible open set, the condition that the $g_c$
generate is also open, and $f_1'(0)\ne0$ is a nonempty open condition.
Inverse reparametrization preserves the quotient and produces graph
form.  For $n=1$, the same assertion is immediate with all $f_a=0$.
Consequently the closed condition
\begin{equation}\label{eq:r5-closed-condition}
                        L_C(\ell)\subseteq\Pol_{n,k}
\end{equation}
holds on the entire curvilinear Quot component, where $\ell$ is the
Pl\"ucker line of a quotient.  More explicitly, compose $L_C$ with
the map $E\to E/(E\cap\Pol_{n,k})$.
Condition \eqref{eq:r5-closed-condition} says that this linear map
vanishes on $\ell$, so it is a closed linear condition on the
Grassmannian.  It includes the points where $L_C(\ell)=0$, and
therefore passes through every specialization in the closure.

\begin{proof}[Proof of Theorem~\ref{thm:r5-products}]
Let $I_c$ be the monomial ideal whose standard monomials are
$\Gamma_c$.  The monomial quotient $\bigoplus_c S/I_c$ has total
length $n$, so Theorem~\ref{thm:r5-bs-input} and
\eqref{eq:r5-closed-condition} apply.
Choose quotient functionals dual to the colored basis
$\{u^\alpha e_c:u^\alpha\in\Gamma_c\}$.
In \eqref{eq:r5-colored-map}, the corresponding column is zero
outside the label block $C_c$.  Thus, up to the ordering sign,
\[
 L_C(\ell)=
 K\left(
   \frac{\prod_{c=1}^r\Delta_{\Gamma_c}(X_{C_c})}
        {\prod_{c=1}^r\prod_{u^\alpha\in\Gamma_c}\alpha!}
     \right).
\]
This is a nonzero line.  It is contained in $\Pol_{n,k}$ by
\eqref{eq:r5-closed-condition}, proving \eqref{eq:r5-product} over
$\C$.  The product, the defining operators, and the Vandermonde all
have rational coefficients.  The assertion is membership in a
finite-dimensional rational linear span, so membership over $\C$
implies membership over $\Q$ and then over every characteristic-zero
field $K$.  Closure under derivatives and polarizations gives the
last assertion.
\end{proof}

Each factor in \eqref{eq:r5-product} is harmonic for its own label
set: a positive diagonal differential operator replaces a monomial
column by a divisor column already present, or by zero.  Since the
label sets are disjoint, the product is harmonic for all $n$ labels.
Theorem~\ref{thm:r5-products} places these products in the operator module.

\subsection{Multiplication on disjoint label sets}

The product theorem extends from lower-diagram alternants to the
whole family of operator modules.  Private row spaces let the
operators act independently on each label block; an invertible row
shear and exact specialization then identify all row spaces.

\begin{theorem}[Disjoint-label multiplication]
\label{thm:r5-disjoint-products}
Let $C_1,\ldots,C_r$ partition $\{1,\ldots,n\}$, with
$|C_c|=n_c>0$.  For arbitrary
$F_c\in\Pol_{n_c,k}$, relabelled on $C_c$, one has
\[
                         \prod_{c=1}^r F_c(X_{C_c})\in\Pol_{n,k}.
\]
These products are associative, commutative, and equivariant under
relabeling and the diagonal row action.  They define canonical
$S_n\times\GL_k$-equivariant maps
\begin{equation}\label{eq:r5-induced-product}
 \operatorname{Ind}_{S_{n_1}\times\cdots\times S_{n_r}}^{S_n}
       \left(\bigotimes_{c=1}^r\Pol_{n_c,k}\right)
                         \longrightarrow\Pol_{n,k}.
\end{equation}
\end{theorem}

\begin{proof}
Introduce $r$ private blocks of $k$ variable rows, with variables
$x_i^{(c,a)}$.  For the $c$th factor take the lower diagram
$\Gamma_c=\{1,u_{c,1},\ldots,u_{c,1}^{n_c-1}\}$ in these $kr$
one-label variables.  Theorem~\ref{thm:r5-products} gives
\[
                    \prod_c\Delta_{C_c}(x^{(c,1)})
                                      \in\Pol_{n,kr}.
\]
Generate $F_c$ using derivatives and polarizations whose source
and target rows both lie in its private block.  Such a global
polarization acts only on the $c$th factor: every other factor is
independent of its source row, and its derivative order is at least
one.  Label-specific partial derivatives have the same property.
By linearity, this produces
$\prod_c F_c(X_{C_c}^{(c)})$ inside $\Pol_{n,kr}$.

Apply the invertible row shear
\[
 x_i^{(c,a)}\longmapsto x_i^{(c,a)}+x_i^{(1,a)}\quad(c>1),
 \qquad x_i^{(1,a)}\longmapsto x_i^{(1,a)}.
\]
Then set all private rows with $c>1$ to zero.
The shear preserves $\Pol_{n,kr}$, and
Corollary~\ref{cor:r3-exact-specialization} sends the result into
$\Pol_{n,k}$.  The resulting polynomial is the desired product
in the common $k$ rows.  The associativity, commutativity, and
equivariance assertions follow from ordinary multiplication in
disjoint polynomial variables.  Frobenius reciprocity gives
\eqref{eq:r5-induced-product}.
\end{proof}

\begin{corollary}[Tensor products under restriction]
\label{cor:r5-product-restriction}
There is a natural injective map of
$\GL_k\times S_{n_1}\times\cdots\times S_{n_r}$-modules
\[
 \bigotimes_c\Pol_{n_c,k}
   \ \lhook\joinrel\longrightarrow\
 \operatorname{Res}_{S_{n_1}\times\cdots\times S_{n_r}}^{S_n}
                                     \Pol_{n,k}.
\]
Consequently, for label partitions $\mu_c\vdash n_c$ and a row
partition $\lambda$ of length at most $k$,
\begin{equation}\label{eq:r5-product-lr}
 \sum_{\nu\vdash n}
       \operatorname{LR}_{\mu_1,\ldots,\mu_r}^{\nu}
                         c_{\lambda,\nu}^{\Pol_{n,k}}
 \ \geq\
 \sum_{\lambda_1,\ldots,\lambda_r}
       \operatorname{LR}_{\lambda_1,\ldots,\lambda_r}^{\lambda}
                     \prod_{c=1}^r
                       c_{\lambda_c,\mu_c}^{\Pol_{n_c,k}}.
\end{equation}
Here $\operatorname{LR}$ denotes an iterated Littlewood--Richardson
coefficient, and the sum on the right is over row partitions of
length at most $k$.
\end{corollary}

\begin{proof}
Multiplication identifies the tensor product of polynomial rings on
the disjoint label sets with the full polynomial ring.  Its
restriction to the tensor product of the indicated subspaces is
therefore injective, and Theorem~\ref{thm:r5-disjoint-products}
places its image in $\Pol_{n,k}$.  Decompose this inclusion under
the row group and the label Young subgroup.  The
Littlewood--Richardson restriction and tensor product rules give
\eqref{eq:r5-product-lr}.
\end{proof}

For fixed label blocks, multiplication is injective as stated in
Corollary~\ref{cor:r5-product-restriction}.  The induced map
\eqref{eq:r5-induced-product} combines all relabelings of those
blocks and can have a kernel.

\subsection{Every positive-weight first layer}
\label{sec:r5-weighted}

The first-occurrence construction works for every positive row
weight.  A minimum-weight column contains all cheaper monomials,
so it remains a lower diagram even when many monomials have equal
weight.  Give the $a$th row a positive integer weight $w_a$.  Write
$P_{n,k}^{(w)}[d]$ for weighted degree $d$ and put
\[
 \begin{aligned}
 W_d(w)&=\Span_K\{u^\alpha:w\cdot\alpha=d\},
 &b_d(w)&=\dim W_d(w),\\
 N_d(w)&=\sum_{e=0}^d b_e(w),
 &N_{-1}(w)&=0.
 \end{aligned}
\]
Order all monomials by nondecreasing weight, breaking ties in any
fixed way.  Let $d_i(w)$ be the weight of the $i$th monomial and set
\[
 \delta_\mu(w)=\sum_i\mu_i d_i(w),\qquad
 \mu^{(d,w)}=
  (\mu_{N_{d-1}(w)+1},\ldots,\mu_{N_d(w)}).
\]
Empty blocks contribute the trivial factor.

\begin{theorem}[Complete generation on every lower weighted face]
\label{thm:r5-weighted-first}
For every positive integral weight vector $w$ and $\mu\vdash n$,
$\delta_\mu(w)$ is the first weighted degree containing $S^\mu$, and
\begin{equation}\label{eq:r5-weighted-equality}
 \Pol_{n,k}^{(w)}[\delta_\mu(w)]_\mu
 =\Harm_{n,k}^{(w)}[\delta_\mu(w)]_\mu
 =P_{n,k}^{(w)}[\delta_\mu(w)]_\mu.
\end{equation}
The multiplicity space of this layer is
\begin{equation}\label{eq:r5-weighted-character}
                  \bigotimes_{d\geq0}
                         \Schur_{\mu^{(d,w)}}(W_d(w)).
\end{equation}
This is an isomorphism for the row torus and for the subgroup of
$\GL_k$ preserving the weighted grading.
\end{theorem}

\begin{proof}
As an $S_n$-module with row action,
$P_{n,k}=(\Sym V)^{\otimes n}$.  Schur--Weyl duality identifies
the $S^\mu$ multiplicity space with $\Schur_\mu(\Sym V)$.
In its semistandard tableau construction, an entry in row $i$ is
at least the $i$th element of the ordered monomial basis.  Its
weight is therefore at least $d_i(w)$.  Filling row $i$ with that
$i$th monomial attains the bound $\delta_\mu(w)$.

Equality in the weight bound requires every box in row $i$ to
contain a monomial of weight exactly $d_i(w)$.  Inside a weight
block the fillings are the semistandard fillings for the partition
$\mu^{(d,w)}$ on $W_d(w)$.  Between different blocks, column
inequalities are automatic.  This proves
\eqref{eq:r5-weighted-character} and the claimed first degree.
Every positive $D_\alpha$ is $S_n$-equivariant and lowers weighted
degree by $w\cdot\alpha>0$.  It annihilates this whole first
isotype, proving equality of the $P$ and $\Harm$ terms in
\eqref{eq:r5-weighted-equality}.

It remains to generate the entire space.  Fix a standard label
tableau $T$ of shape $\mu$, with label sets $C_c$ in its columns.
Let $F$ be any semistandard monomial filling attaining
$\delta_\mu(w)$, and let $\Gamma_c$ be the monomial set in its
$c$th column.  Column strictness makes these monomials distinct.
For each weight strictly below the last weight in that column,
the column crosses the complete corresponding block of rows.
It must therefore contain every monomial of that weight; in the
last block it contains a subset.  Since all $w_a$ are positive,
a proper monomial divisor has strictly smaller weight.  Thus
every $\Gamma_c$ is a lower diagram.

Let $m_F$ be the monomial in $P_{n,k}$ obtained by putting the
monomial in each box of $F$ at the label in the same box of $T$.
If $b_T$ is the column antisymmetrizer, then
\[
                     b_Tm_F
                       =\prod_c\Delta_{\Gamma_c}(X_{C_c})
\]
up to the fixed column ordering signs.  Theorem~\ref{thm:r5-products}
places this polynomial in $\Pol_{n,k}$.  Applying the row
symmetrizer $a_T$ of the label tableau also preserves $\Pol_{n,k}$.
The standard
semistandard construction says that the vectors $a_Tb_Tm_F$, as
$F$ ranges over these minimum-weight fillings, span the
first-weight part of the image of the Young symmetrizer $a_Tb_T$.
That image contains one fixed Specht line tensored with the whole
multiplicity space \eqref{eq:r5-weighted-character}.
Its $S_n$-span is the complete $\mu$-isotype.  Since $\Pol_{n,k}$
is $S_n$-stable, this proves the remaining containment in
\eqref{eq:r5-weighted-equality}.
\end{proof}

\begin{proof}[Proof of Theorem~\ref{thm:r5-complete-first}]
Take $w=(1,\ldots,1)$ in Theorem~\ref{thm:r5-weighted-first}.
Then $W_j(w)=\Sym^jV$, $b_j(w)=b_j$, and $N_j(w)=N_j$.
\end{proof}

The proof generates monomial fillings one by one before assembling
their row representations.  It therefore carries every irreducible
constituent and every plethystic multiplicity in
\eqref{eq:r5-weighted-character} into the operator module.

\begin{corollary}[The whole positive-weight lower boundary]
\label{cor:r5-lower-boundary}
For $\mathcal M=\Pol,\Harm$, let
$\operatorname{ch}_\mu^{\mathcal M}(\bq)$ be the row-torus character
of the $S^\mu$ multiplicity space in $\mathcal M_{n,k}$.  For every
$w\in\N_{>0}^k$, the terms minimizing $w\cdot\beta$ in the two
characters are identical, including their coefficients.  They are
\begin{equation}\label{eq:r5-face-character}
 \prod_{d\geq0}
 s_{\mu^{(d,w)}}
     \left[\sum_{w\cdot\alpha=d}q^\alpha\right],
\end{equation}
and their common minimum is $\delta_\mu(w)$.
Consequently the generation defect has no component on any face
of the $\mu$-character support exposed by a strictly positive
weight vector.
\end{corollary}

\begin{proof}
Theorem~\ref{thm:r5-weighted-first} is equality of the corresponding
row-torus modules.  The character of a Schur functor of $W_d(w)$
is the plethystic factor in \eqref{eq:r5-face-character}.
Every face exposed by a strictly positive real vector is also
exposed by a strictly positive rational vector: the relevant
relatively open normal cone is defined by rational linear
equalities and strict inequalities.  Clearing denominators
reduces that face to the integral-weight case.
\end{proof}

Strict positivity makes every proper monomial divisor cheaper and
every positive invariant differential operator degree lowering.
These are the two properties that drive the weighted proof.

